% ============================================================================
% 张博 — 个人简历（中文版）
% 编译：xelatex cv_zh.tex（两次）
% ============================================================================
\documentclass[10pt]{article}
\usepackage{bo-cv}

\cvfooter{}
\fancyfoot[R]{}

\newcommand{\cvcontrib}[1]{%
  {\color{muted}\fontsize{9.5}{12}\selectfont 个人贡献：#1}\par\vspace{0.05em}%
}

\begin{document}

% ----- 标题 ----------------------------------------------------------------
\cvheader[image.png]%
  {张博}%
  {}%
  {本科生 \cvsep\ 清华大学自动化系}%
  {%
    \cvicon{\faEnvelope[regular]}\href{mailto:bo-zhang23@mails.tsinghua.edu.cn}{bo-zhang23@mails.tsinghua.edu.cn}\quad\cvicon{\faGlobe}\href{https://bozhang.site}{bozhang.site}\\%
    \cvicon{\faGithub}\href{mailto:226653803@qq.com}{226653803@qq.com}\quad\cvicon{\faPhone}\href{tel:+8615109515899}{15109515899}%
  }

% ----- 教育背景 ------------------------------------------------------------
\section{教育背景}

\cventry%
  {清华大学 \cvsep\ 自动化系}%
  {工学学士，自动化 \cvsep\ 多模态智能因材施教方向班}%
  {2023.09 -- 2027.07（预计）}%
  {%
    必限学分绩 3.5/4.0 \cvsep\ 专业内排名 61/172 \\
    {\itshape 主要修读课程}：模式识别与机器学习 \cvsep\ 生成式人工智能 \cvsep\ 运筹学 \cvsep\ 计算机视觉 \cvsep\ 运筹学 \cvsep\ 知识工程%
  }

% ----- 发表论文 ------------------------------------------------------------
\section{发表论文}

\cvtheme{AI for Science}{具身 LLM 智能体与科学软件交互框架}

\cvpub%
  {Grounding {LLMs} in Scientific Discovery via Embodied Actions}%
  {\textbf{Bo~Zhang}*, Jinfeng~Zhou*, Yuxuan~Chen, Jianing~Yin, Minlie~Huang, Hongning~Wang$^\dagger$}%
  {ICML 2026 \cvsep\ {*} 共同第一作者 \cvsep\ 隶属：清华大学}%
  {}
\cvcontrib{构建科学软件实时观测-动作交互框架，支持智能体执行、反馈获取与实验验证。}

\cvtheme{计算机代理智能体}{精细视觉定位与免训练推理增强}

\cvpub%
  {{BAMI}: Training-Free Bias Mitigation in {GUI} Grounding}%
  {\textbf{Bo~Zhang}*, Borui~Zhang*, Bo~Wang*, Wenzhao~Zheng, Yuhao~Cheng, Liang~Tang, Yiqiang~Yan, Jie~Zhou, Jiwen~Lu$^\dagger$}%
  {CVPR 2026 \cvsep\ {*} 共同第一作者 \cvsep\ 隶属：清华大学 / 联想研究院}%
  {}
\cvcontrib{设计偏置归因分析流程与免训练推理增强策略，并完成 GUI grounding 实验验证。}

% ----- 研究与项目 ----------------------------------------------------------
\section{研究与项目}

\cvtheme{计算机代理智能体}{框架可审计性与动作-观测空间完备性}

\cventry%
  {SRT --- GUI 桌面视觉理解}%
  {清华大学本科生研究训练计划（SRT） \cvsep\ 科研训练}%
  {2024 -- 2025}%
  {构建桌面 GUI 视觉理解数据与评测流程，分析 GUI grounding 空间偏置，为 BAMI 方法设计提供实验依据。}

\cventry%
  {Syll --- 自部署多模态智能体}%
  {核心贡献者 \cvsep\ 开源项目 \cvsep\ \href{https://github.com/THU-SAGE/syll}{THU-SAGE/syll}}%
  {2026 -- 至今}%
  {设计自部署多通道桌面智能体执行框架，统一 GUI / CLI / API 动作空间；以截图、日志与工作流录制-回放支持过程审计。}

\cvtheme{教育智能体}{多智能体平台与课程级 AI 助教}

\cventry%
  {AICosmos --- 教育多智能体平台}%
  {贡献者 \cvsep\ 清华大学计算机系孵化项目}%
  {2025 -- 至今}%
  {参与多智能体教学流程设计与平台功能实现，支持情景化教学、多模态交互；\mbox{已上线} \href{https://aicosmos.ai}{aicosmos.ai}。}

\cventry%
  {AEI4U --- 模拟电子技术 AI 助教}%
  {项目负责人 \cvsep\ 基于 Dify-THU 平台}%
  {2025}%
  {课程级 AI 助教，集成电路拓扑识别、SPICE 仿真接入与个性化学习画像；获清华大学第 43 届\mbox{挑战杯} \mbox{\textbf{特等奖}}。}

% ----- 荣誉、专利与学生工作 ------------------------------------------------
\section{荣誉、专利与学生工作}

\cvrow{荣誉}{清华大学第 43 届\mbox{挑战杯}\mbox{\textbf{特等奖}}（2025）}
\cvrow{专利}{%
  《基于具身科学代理的科学发现检验方法及系统》\newline%
  {\color{muted}\fontsize{8.7}{11}\selectfont 申请号：202610216049.3 \cvsep\ 发明人：王宏宁、周金锋、\textbf{张博} \cvsep\ 申请人：清华大学 \cvsep\ 校内终审通过}\newline%
  《一种基于示教-推理-执行闭环的多模态智能体系统》\newline%
  {\color{muted}\fontsize{8.7}{11}\selectfont 申请号：202610548720.4 \cvsep\ 发明人：鲁继文、周杰、\textbf{张博}、张博睿、江承昊 \cvsep\ 申请人：清华大学 \cvsep\ 校内终审通过}%
}
\cvrow{学生工作}{多模态智能因材施教方向班 班长（2024 -- 至今）\cvsep\ 自动化系篮球队 队长}

% ----- 技能与语言 ----------------------------------------------------------
\section{技能与语言}

\cvrow{研究方法}{智能体框架搭建 \cvsep\ 免训练推理增强 \cvsep\ 偏置归因分析 }
\cvrow{工程能力}{Python \cvsep\ PyTorch \cvsep\ FastAPI \cvsep\ LiteLLM \cvsep\ Docker \cvsep\ MATLAB \cvsep\ Dify}
\cvrow{语言}{普通话（母语）\cvsep\ 英语（可进行学术阅读、写作与口头交流）}

\end{document}
